\documentclass[11pt,a4paper]{article}

\usepackage[T2A]{fontenc}
\usepackage[utf8]{inputenc}
\usepackage[russian]{babel}
\usepackage[a4paper,hmargin=2cm,vmargin=1.7cm]{geometry}
\usepackage{amsmath,amssymb}
\usepackage{array,booktabs}
\usepackage{enumitem}

\setlength{\parindent}{0pt}
\setlength{\parskip}{4pt}
\setlist{nosep,leftmargin=*}

% Заголовок примера
\newcommand{\example}[1]{\par\medskip\noindent\textbf{#1}\par\nopagebreak\smallskip}

\begin{document}

\begin{center}
{\Large\bfseries Линейная зависимость функций. \\[2pt]  Формула Лиувилля "--- Остроградского}
\end{center}

% =====================================================================
\section*{Линейная зависимость функций}

\subsection*{Определение}

Функции $y_1(x), \dots, y_n(x)$ \textbf{линейно зависимы} на интервале $(a, b)$, если найдутся числа $\alpha_1, \dots, \alpha_n$, \textbf{не все равные нулю}, такие что
\[
  \alpha_1 y_1(x) + \alpha_2 y_2(x) + \dots + \alpha_n y_n(x) \equiv 0 \quad \text{на } (a, b).
\]
Если тождество выполняется только при $\alpha_1 = \dots = \alpha_n = 0$, функции \textbf{линейно независимы}.

Иначе говоря, функции зависимы, если одну из них можно выразить через остальные.

\subsection*{Случай двух функций}

Две функции $y_1$, $y_2$ (при $y_2 \ne 0$) линейно зависимы тогда и только тогда, когда их отношение постоянно:
\[
  \frac{y_1}{y_2} \equiv \text{const}
  \quad\Leftrightarrow\quad
  \text{зависимы}\,, \qquad
  \frac{y_1}{y_2} \not\equiv \text{const}
  \quad\Leftrightarrow\quad
  \text{независимы}
\]

\example{Пример 1 (Филиппов № 642, 643)}

Являются ли линейно зависимыми функции:
\quad а) $6x + 9$, \ $8x + 12$; \qquad б) $\sin x$, \ $\cos x$?

Решение:

а) Выносим общие множители:
\[
  \frac{6x + 9}{8x + 12} = \frac{3(2x + 3)}{4(2x + 3)} = \frac34 = \text{const}
\]
Отношение постоянно, значит, функции линейно зависимы: $4(6x + 9) - 3(8x + 12) \equiv 0$.

б) Отношение
\[
  \frac{\sin x}{\cos x} = \operatorname{tg} x \ne \text{const}
\]
Отношение меняется, значит, функции линейно независимы.

Ответ: а) зависимы; б) независимы.

% =====================================================================
\section*{Матрица Вронского. Определитель Вронского (вронскиан)}

\subsection*{Определение}

\textbf{Матрица Вронского}, для $n$ функций, имеющих производные до порядка $n - 1$:
\[
  W(x) =
  \begin{pmatrix}
    y_1 & y_2 & \dots & y_n \\
    y_1' & y_2' & \dots & y_n' \\
    \vdots & \vdots & & \vdots \\
    y_1^{(n-1)} & y_2^{(n-1)} & \dots & y_n^{(n-1)}
  \end{pmatrix}
\]
В первой строке "--- сами функции, в каждой следующей "--- производные предыдущей строки.

\textbf{Определитель Вронского (вронскиан):}
\[
  w(x) = \det W(x)
\]


\subsection*{Связь с линейной зависимостью}

\textbf{Теорема.} Если функции $y_1, \dots, y_n$ линейно зависимы на $(a, b)$, то $w(x) \equiv 0$ на $(a, b)$.

\textbf{Как применять:} если $w(x_0) \ne 0$ хотя бы в одной точке, то функции \textbf{линейно независимы}.

Обратное для произвольных функций неверно: из $w(x) \equiv 0$ зависимость не следует. Для решений одного линейного уравнения обратное верно "--- см. следующий раздел.

\example{Пример 2 (Филиппов № 649)}

Являются ли линейно зависимыми функции $x$, \ $e^{x}$, \ $x e^{x}$?

Решение:

Функций три, поэтому в определителе три строки: функции, первые и вторые производные.
\[
  (x e^{x})' = (x + 1)e^{x}\,, \qquad (x e^{x})'' = (x + 2)e^{x}
\]
\[
  w(x) =
  \begin{vmatrix}
    x & e^{x} & x e^{x} \\
    1 & e^{x} & (x + 1)e^{x} \\
    0 & e^{x} & (x + 2)e^{x}
  \end{vmatrix}
\]

Выносим $e^{x}$ из второго и третьего столбцов:
\[
  w(x) = e^{2x}
  \begin{vmatrix}
    x & 1 & x \\
    1 & 1 & x + 1 \\
    0 & 1 & x + 2
  \end{vmatrix}
\]

Раскладываем по первому столбцу:
\[
  w(x) = e^{2x}\Bigl(x \cdot \bigl((x + 2) - (x + 1)\bigr) - 1 \cdot \bigl((x + 2) - x\bigr)\Bigr)
  = e^{2x}(x - 2)
\]

$w(x) \not\equiv 0$: например, $w(0) = -2 \ne 0$. Значит, функции линейно независимы.


Ответ: функции линейно независимы.

% =====================================================================
\section*{Фундаментальная система решений}

\subsection*{Вронскиан решений}

Рассмотрим линейное однородное уравнение с непрерывными на $(a, b)$ коэффициентами:
\[
  y'' + p(x)\, y' + q(x)\, y = 0
\]
(для уравнения порядка $n$ всё аналогично).

\textbf{Теорема.} Если $y_1$, $y_2$ "--- решения этого уравнения, то возможны только два случая:
\begin{itemize}
  \item $w(x) \equiv 0$ на $(a, b)$ "--- решения линейно зависимы;
  \item $w(x) \ne 0$ во всех точках $(a, b)$ "--- решения линейно независимы.
\end{itemize}
Поэтому для решений достаточно проверить $w(x)$ в одной точке.

\subsection*{ФСР и общее решение}

\textbf{Фундаментальная система решений (ФСР)} уравнения порядка $n$ "--- это $n$ линейно независимых решений. Для уравнения второго порядка "--- два решения $y_1$, $y_2$ с $w(x) \ne 0$.

Общее решение "--- их линейная комбинация:
\[
  y = C_1 y_1 + C_2 y_2
\]

На занятии 1 в задаче Коши мы решали систему с матрицей Вронского $W(x_0)$. Она обратима как раз потому, что $w(x_0)=\det W(x_0) \ne 0$ для ФСР.

\subsection*{Алгоритм: проверить, что функции образуют ФСР}

\textbf{1.} Подставить каждую функцию в дифференциальное уравнение и убедиться, что это решение.

\textbf{2.} Число функций равно порядку уравнения.

\textbf{3.} Вычислить $w(x)$ и убедиться, что $w(x) \ne 0$.

\textbf{4.} Записать общее решение $y = C_1 y_1 + C_2 y_2$.

\example{Пример 3}

Проверить, что функции $y_1 = x$, $y_2 = x^2 - 1$ образуют ФСР уравнения
\[
  (x^2 + 1)\, y'' - 2x\, y' + 2y = 0\,,
\]
и записать общее решение.

Решение:

\textbf{1.} Подстановка.

$y_1 = x$: \ $y_1' = 1$, $y_1'' = 0$:
\[
  (x^2 + 1) \cdot 0 - 2x \cdot 1 + 2x = 0
\]

$y_2 = x^2 - 1$: \ $y_2' = 2x$, $y_2'' = 2$:
\[
  (x^2 + 1) \cdot 2 - 2x \cdot 2x + 2(x^2 - 1) = 2x^2 + 2 - 4x^2 + 2x^2 - 2 = 0
\]
Обе функции "--- решения.

\textbf{2.} Функций две, порядок уравнения 2.

\textbf{3.} Вронскиан:
\[
  w(x) =
  \begin{vmatrix} x & x^2 - 1 \\ 1 & 2x \end{vmatrix}
  = 2x^2 - (x^2 - 1) = x^2 + 1 \ne 0
\]
Решения линейно независимы, это ФСР.

\textbf{4.} Общее решение "--- линейная комбинация решений ФСР.

Ответ:
\[
  y = C_1 x + C_2 (x^2 - 1)
\]

% =====================================================================
\section*{Формула Лиувилля "--- Остроградского}


Для любых двух решений уравнения $y'' + p(x)\, y' + q(x)\, y = 0$:
\[
  w(x) = C e^{-\int p(x)\,dx}
\]
Вронскиан находится \textbf{без знания решений}, по одному коэффициенту $p(x)$.

\textbf{Важно:} $p(x)$ "--- коэффициент при $y'$, когда коэффициент при $y''$ равен $1$. Для уравнения $a_0(x)\, y'' + a_1(x)\, y' + a_2(x)\, y = 0$ сначала делим на $a_0$:
\[
  p(x) = \frac{a_1(x)}{a_0(x)}
\]
Для уравнения порядка $n$ формула та же, $p(x)$ "--- коэффициент при $y^{(n-1)}$.

\subsection*{Алгоритм: второе решение по известному $y_1$ и общее решение для ДУ 2-го порядка}

Если одно решение $y_1$ известно, формула Лиувилля даёт второе.

\textbf{1.} Разделить уравнение на коэффициент при $y''$ и выписать $p(x)$.

\textbf{2.} Найти вронскиан по формуле Лиувилля "--- Остроградского, взяв $C = 1$ (нам нужно какое-нибудь одно второе решение):
\[
  w(x) = e^{-\int p(x)\,dx}
\]

\textbf{3.} Второе решение вычисляется по формуле:
\[
  y_2 = y_1 \int \frac{w}{y_1^2}\,dx
\]
Числовой множитель перед $y_2$ можно отбросить: он войдёт в $C_2$.

\textbf{4.} Записать общее решение $y = C_1 y_1 + C_2 y_2$.

\example{Пример 4 (Филиппов № 683)}

Найти общее решение уравнения
\[
  x y'' - (2x + 1)\, y' + (x + 1)\, y = 0\,, \qquad x > 0\,,
\]
если известно частное решение $y_1 = e^{x}$ (в книге $y_1$ не дано, его нужно подобрать; здесь оно дано нами).

Решение:

\textbf{1.} Делим на $x$:
\[
  y'' - \frac{2x + 1}{x}\, y' + \frac{x + 1}{x}\, y = 0\,, \qquad
  p(x) = -\frac{2x + 1}{x} = -2 - \frac1x
\]

\textbf{2.} Вронскиан:
\[
  -\int p\,dx = \int \left(2 + \frac1x\right) dx = 2x + \ln x\,, \qquad
  w(x) = e^{-\int p(x)\,dx} = e^{2x + \ln x} = x e^{2x}
\]

\textbf{3.} Второе решение:
\[
  y_2 = y_1 \int \frac{w}{y_1^2}\,dx = e^x\int\frac{xe^{2x}}{e^{2x}}\,dx = e^x \int x\,dx = \frac{x^2}{2}\, e^{x}
\]
Множитель $\frac12$ отбрасываем: $y_2 = x^2 e^{x}$.

\textbf{4.} Общее решение $y = C_1 y_1 + C_2 y_2$.

Ответ:
\[
  y = (C_1 + C_2 x^2)\, e^{x}
\]

\example{Пример 5 (Филиппов № 684)}

Найти общее решение уравнения
\[
  x y'' + 2y' - x y = 0\,,
\]
если известно частное решение $y_1 = \dfrac{e^{x}}{x}$.

Решение:

\textbf{1.} Делим на $x$:
\[
  y'' + \frac{2}{x}\, y' - y = 0\,, \qquad p(x) = \frac{2}{x}
\]

\textbf{2.} Вронскиан:
\[
  -\int p\,dx = -2\ln|x|\,, \qquad
  w = e^{-\int p(x)\,dx} = e^{-2\ln|x|} = \frac{1}{x^2}
\]

\textbf{3.} Второе решение:
\[
  y_2 = y_1 \int \frac{w}{y_1^2}\,dx = \dfrac{e^{x}}{x} \int e^{-2x}\,dx= -\frac12\, e^{-2x} \cdot \frac{e^{x}}{x} = -\frac12 \cdot \frac{e^{-x}}{x}
\]
Множитель $-\frac12$ отбрасываем: $y_2 = \dfrac{e^{-x}}{x}$.

\textbf{4.} Общее решение $y = C_1 y_1 + C_2 y_2$.

Ответ:
\[
  y = \frac{C_1 e^{x} + C_2 e^{-x}}{x}
\]

% =====================================================================

\section*{Домашнее задание}

\subsection*{Линейная зависимость и вронскиан}
\begin{itemize}[topsep=2pt]
  \item \textbf{Филиппов № 641} "--- являются ли линейно зависимыми
  \begin{itemize}
    \item функции: $x + 2$, \ $x - 2$
    \item ответ: линейно независимы
  \end{itemize}
  \item \textbf{Филиппов № 647} "--- являются ли линейно зависимыми; если да, найти коэффициенты
  \begin{itemize}
    \item функции: $x^2 - x + 3$, \ $2x^2 + x$, \ $2x - 4$
    \item ответ: линейно зависимы, $4(x^2 - x + 3) - 2(2x^2 + x) + 3(2x - 4) \equiv 0$
    \item в книге ответ только «да», коэффициенты добавлены нами
  \end{itemize}
  \item \textbf{Краснов № 392} "--- найти определитель Вронского
  \begin{itemize}
    \item функции: $e^{-x}$, \ $x e^{-x}$
    \item ответ: $w = e^{-2x}$
  \end{itemize}
\end{itemize}

\subsection*{Формула Лиувилля "--- Остроградского: общее решение по частному}
\begin{itemize}[topsep=2pt]
  \item \textbf{Филиппов № 692}
  \begin{itemize}
    \item уравнение: $y'' + 4xy' + (4x^2 + 2)\, y = 0$, \ $y_1 = e^{-x^2}$
    \item ответ: $y = (C_1 + C_2 x)\, e^{-x^2}$
    \item в книге дано $y_1 = e^{ax^2}$, число $a$ нужно найти; здесь $a = -1$ подставлено нами
  \end{itemize}
  \item \textbf{Филиппов № 685}
  \begin{itemize}
    \item уравнение: $y'' - 2(1 + \operatorname{tg}^2 x)\, y = 0$, \ $y_1 = \operatorname{tg} x$
    \item ответ: $y = C_1 \operatorname{tg} x + C_2 (1 + x \operatorname{tg} x)$
  \end{itemize}
  \item \textbf{Филиппов № 693} "--- нужно интегрирование по частям
  \begin{itemize}
    \item уравнение: $x y'' - (2x + 1)\, y' + 2y = 0$, \ $y_1 = e^{2x}$
    \item ответ: $y = C_1 (2x + 1) + C_2 e^{2x}$
    \item в книге $y_1$ не дано, его нужно подобрать; здесь оно дано нами
  \end{itemize}
\end{itemize}

% =====================================================================
\newpage

\section*{Дополнительные примеры}

\example{Пример Д1 (Филиппов № 650, зависимость через тождество)}

Являются ли линейно зависимыми функции $1$, \ $\sin^2 x$, \ $\cos 2x$?

Решение:

Формула двойного угла: $\cos 2x = 1 - 2\sin^2 x$. Переносим всё в одну сторону:
\[
  1 \cdot 1 - 2 \cdot \sin^2 x - 1 \cdot \cos 2x \equiv 0
\]
Коэффициенты $1, -2, -1$ не равны нулю, значит, функции линейно зависимы. По теореме $w \equiv 0$ (это можно проверить и прямым вычислением).

Ответ: функции линейно зависимы.

\textbf{Совет:} если в наборе есть тригонометрические, гиперболические функции или логарифмы, сначала ищите тождество, связывающее их.

\example{Пример Д2 (Филиппов № 660, $w \equiv 0$, но функции независимы)}

Являются ли линейно зависимыми функции $y_1 = x^2$, \ $y_2 = x|x|$ на интервале $(-1, 1)$?

Решение:

Раскрываем модуль:
\[
  y_2 = x|x| =
  \begin{cases}
    \phantom{-}x^2, & x \ge 0\,,\\
    -x^2, & x < 0\,.
  \end{cases}
\]
Отношение $\dfrac{y_2}{y_1}$ равно $1$ при $x > 0$ и $-1$ при $x < 0$ "--- оно не постоянно на $(-1, 1)$. Значит, функции линейно независимы.

Вронскиан. Так как $(x|x|)' = 2|x|$:
\[
  w = x^2 \cdot 2|x| - 2x \cdot x|x| = 2x^2|x| - 2x^2|x| = 0 \quad\text{при всех } x
\]

Противоречия с теоремой о вронскиане решений нет: эти функции не могут быть решениями одного уравнения $y'' + p(x)y' + q(x)y = 0$ с непрерывными на $(-1, 1)$ коэффициентами.

Ответ: функции линейно независимы, хотя $w \equiv 0$.

\example{Пример Д3 (почему таблица ФСР из занятия 1 даёт независимые функции)}

Вычислим вронскиан для пар функций из таблицы ФСР занятия 1.

Разные вещественные корни $\lambda_1 \ne \lambda_2$:
\[
  w[e^{\lambda_1 x}, e^{\lambda_2 x}] =
  \begin{vmatrix} e^{\lambda_1 x} & e^{\lambda_2 x} \\ \lambda_1 e^{\lambda_1 x} & \lambda_2 e^{\lambda_2 x} \end{vmatrix}
  = (\lambda_2 - \lambda_1)\, e^{(\lambda_1 + \lambda_2)x} \ne 0
\]

Кратный корень $\lambda$:
\[
  w[e^{\lambda x}, x e^{\lambda x}] =
  \begin{vmatrix} e^{\lambda x} & x e^{\lambda x} \\ \lambda e^{\lambda x} & (1 + \lambda x) e^{\lambda x} \end{vmatrix}
  = e^{2\lambda x}\bigl((1 + \lambda x) - \lambda x\bigr) = e^{2\lambda x} \ne 0
\]

Комплексные корни $\alpha \pm i\beta$, $\beta \ne 0$. Производные:
\[
  (e^{\alpha x}\cos\beta x)' = e^{\alpha x}(\alpha\cos\beta x - \beta\sin\beta x)\,, \qquad
  (e^{\alpha x}\sin\beta x)' = e^{\alpha x}(\alpha\sin\beta x + \beta\cos\beta x)
\]
\[
\begin{aligned}
  w &= e^{2\alpha x}\bigl(\cos\beta x\,(\alpha\sin\beta x + \beta\cos\beta x) - \sin\beta x\,(\alpha\cos\beta x - \beta\sin\beta x)\bigr)\\
    &= \beta e^{2\alpha x}(\cos^2\beta x + \sin^2\beta x) = \beta e^{2\alpha x} \ne 0
\end{aligned}
\]

Во всех трёх случаях $w \ne 0$, поэтому пары из таблицы действительно образуют ФСР.

\example{Пример Д4 (Филиппов № 686, формула Лиувилля)}

Найти общее решение уравнения
\[
  x(x - 1)\, y'' - x y' + y = 0\,,
\]
если известно частное решение $y_1 = x$ (в книге $y_1$ не дано, его нужно подобрать; здесь оно дано нами).

Решение:

\textbf{1.} Делим на $x(x - 1)$ (рассматриваем интервал, где $x \ne 0$, $x \ne 1$):
\[
  p(x) = \frac{-x}{x(x - 1)} = -\frac{1}{x - 1}
\]

\textbf{2.} Вронскиан:
\[
  -\int p\,dx = \int \frac{dx}{x - 1} = \ln|x - 1|\,, \qquad
  w = |x - 1|
\]
Знак модуля войдёт в константу, берём $w = x - 1$.

\textbf{3.} Второе решение:
\[
  y_2 = y_1 \int \frac{w}{y_1^2}\,dx = x \int \frac{x - 1}{x^2}\,dx
  = x \int \left(\frac1x - \frac{1}{x^2}\right) dx
  = x\left(\ln|x| + \frac1x\right) = x\ln|x| + 1
\]

\textbf{4.} Общее решение $y = C_1 y_1 + C_2 y_2$.

Ответ:
\[
  y = C_1 x + C_2 (x\ln|x| + 1)
\]

\example{Пример Д5 (Филиппов № 697, подбор первого решения)}

Найти общее решение уравнения
\[
  (x^2 + 1)\, y'' - 2y = 0
\]

Решение:

Частное решение не дано "--- сначала подбираем его в виде многочлена, затем применяем алгоритм.

\textbf{Подбор $y_1$: степень многочлена.} Подставляем $y = x^n + \dots$ и оставляем только старшие степени $x$:
\[
  x^2 \cdot n(n - 1)x^{n-2} - 2x^n = \bigl(n(n - 1) - 2\bigr)x^n = 0
  \quad\Rightarrow\quad
  n^2 - n - 2 = 0\,, \quad n = 2 \text{ или } n = -1
\]
Степень многочлена "--- натуральное число, поэтому $n = 2$.

\textbf{Подбор $y_1$: коэффициенты.} Ищем $y_1 = x^2 + ax + b$, \ $y_1'' = 2$:
\[
  (x^2 + 1) \cdot 2 - 2(x^2 + ax + b) = -2ax + (2 - 2b) = 0
  \quad\Rightarrow\quad
  a = 0\,, \quad b = 1
\]
Итак, $y_1 = x^2 + 1$.

\textbf{1.} Делим на $x^2 + 1$:
\[
  y'' - \frac{2}{x^2 + 1}\, y = 0\,, \qquad p(x) = 0
\]

\textbf{2.} Вронскиан: $w = e^{0} = 1$.

\textbf{3.} Второе решение:
\[
  y_2 = y_1 \int \frac{w}{y_1^2}\,dx = (x^2 + 1) \int \frac{dx}{(x^2 + 1)^2}
\]

Интеграл берём заменой $x = \operatorname{tg} t$: \ $dx = \dfrac{dt}{\cos^2 t}$, \ $x^2 + 1 = \dfrac{1}{\cos^2 t}$:
\[
  \int \frac{dx}{(x^2 + 1)^2} = \int \cos^2 t\,dt = \frac{t}{2} + \frac{\sin 2t}{4}
  = \frac12 \operatorname{arctg} x + \frac{x}{2(x^2 + 1)}
\]
Здесь $\sin 2t = \dfrac{2\operatorname{tg} t}{1 + \operatorname{tg}^2 t} = \dfrac{2x}{1 + x^2}$.

Умножаем на $y_1 = x^2 + 1$ и отбрасываем множитель $\frac12$:
\[
  y_2 = x + (x^2 + 1)\operatorname{arctg} x
\]

\textbf{4.} Общее решение $y = C_1 y_1 + C_2 y_2$.

Ответ:
\[
  y = C_1 (x^2 + 1) + C_2 \bigl(x + (x^2 + 1)\operatorname{arctg} x\bigr)
\]

\end{document}
